\section{章末题与参考分析的逐题导航}
\label{sec:appG-map}

章末题使用“章号-题号”，G1--G18是另设的补充题集。下表逐题列出相关分析及作答范围；“开放”表示不给唯一标准答案，所列入口与要点帮助检验论证，不能把相关G题当作原题的完整答案。计算行给出给定条件下的结果，仍须说明假设。点击章末题号可回到原题，点击 G 题号可直达相应的展开分析。

\begin{small}
\setlength{\tabcolsep}{3pt}
\begin{longtable}{@{}p{1.1cm}p{4.1cm}p{10.1cm}@{}}
\toprule
题号 & 相关分析入口 & 作答要点与答案范围 \\
\midrule\endfirsthead
\toprule 题号 & 相关分析入口 & 作答要点与答案范围 \\\midrule\endhead
\bottomrule\endfoot

\hyperref[ex:ch01-1]{1-1} & \hyperref[sol:g01]{G1}/\hyperref[sol:g02]{G2}；式\eqref{eq:ak} & 开放判定：把观察到的改变落到具体组件，核对版本和后续加载。只在本次任务里使用的反思可随轨迹归档。 \\
\hyperref[ex:ch01-2]{1-2} & 表\ref{tab:harness}；附录\ref{app:stats} & 开放设计：可描述点估计排序，不能只凭均值区分效应与波动；先指定比较，再补逐题差、重复、选择规则与差异区间。 \\
\hyperref[ex:ch01-3]{1-3} & \hyperref[sol:g04]{G4}；公理1 & 开放设计：评价可变时另保留外部参照，并报告规则变化对分数可比性的影响。 \\
\hyperref[ex:ch01-4]{1-4} & \hyperref[sol:g06]{G6}；第\ref{ch03}章 & 开放设计：固定起点与完整资源，对齐候选选择；把任务内重试同跨任务继承分开。 \\
\hyperref[ex:ch02-1]{2-1} & 表\ref{tab:roles}；式\eqref{eq:delta}；\hyperref[sol:g04]{G4} & 开放判定：列出修改器可编辑的对象。若任务或评分随版本而变，差值仍可计算，但评价系统收益还需共同的外部参照。 \\
\hyperref[ex:ch02-2]{2-2} & 第\ref{sec:ch02-states}节；\hyperref[sol:g05]{G5} & 开放设计：预定任务序列和重复次数，比较逐次复位与共用目录两组。文件检查核对起点，题级通过率差及区间估计效果；即使点差为零，也应报告区间。 \\
\hyperref[ex:ch02-3]{2-3} & 第\ref{sec:ch02-gate}节；式\eqref{eq:gate} & 开放判定：隔离外部写入后，影子版本可并行接收输入，输出不交付；灰度版本会承担部分正式流量。二元门缺少发布前的影子观察阶段。 \\
\hyperref[ex:ch02-4]{2-4} & 表\ref{tab:states}；第\ref{sec:ch02-reset-matrix}节 & 开放设计：按新题清单和选定版本核对起始文件，再查缓存或外部状态。用持久组件的变更与加载记录判断哪里发生了更新。 \\
\hyperref[ex:ch03-1]{3-1} & \hyperref[sol:g06]{G6}/\hyperref[sol:g07]{G7}；表\ref{tab:ch03-verdict} & 开放设计：起始经验、选择机会、完整费用与任务覆盖均可能混淆，逐项设置匹配对照。 \\
\hyperref[ex:ch03-2]{3-2} & 式\eqref{eq:samplesize}、\eqref{eq:ch03-power}；\hyperref[sol:g05]{G5} & 计算：用题设标准差0.29，半宽近似需359题，双侧功效近似需734题，均向上取整。两者都假设独立配对题、稳定方差；后者还用到备择效应和目标功效。 \\
\hyperref[ex:ch03-3]{3-3} & \hyperref[sol:g06]{G6}；第\ref{ch03}章题集划分 & 开放设计：选择过程和独立终评缺一不可；额外候选和封存评价都进入预算。 \\
\hyperref[ex:ch03-4]{3-4} & 第\ref{sec:ch03-control}节；表\ref{tab:ch15-reader} & 开放判定：以同协议的目标主干静态组作对照。格式偏差与答题分可并列观察，两项汇总量无法算出格式偏差造成的失分比例。 \\
\hyperref[ex:ch03-5]{3-5} & 第\ref{ch03}章；\hyperref[sol:g10]{G10} & 开放设计：族内扩题改善条件内估计；新族扩大覆盖。按目标总体决定权重和抽样单位。 \\
\hyperref[ex:ch04-1]{4-1} & 表\ref{tab:ch04-six}；\hyperref[sol:g15]{G15} & 开放设计：列观测来源、可重复条件和不可重放的外部状态，允许多个合理方案。 \\
\hyperref[ex:ch04-2]{4-2} & 第\ref{ch04}章标签与迁移；\hyperref[sol:g07]{G7} & 开放构造：正确但不适用的检索可损害；有噪声的材料也可能保留有用信息，须给具体路径。 \\
\hyperref[ex:ch04-3]{4-3} & \hyperref[sol:g08]{G8}；表\ref{tab:ch04-adj} & 开放设计：臂含义、额外预算、选择规则及反馈来源须先核对，点估计反序不是层级判决。 \\
\hyperref[ex:ch04-4]{4-4} & \hyperref[sol:g05]{G5}/\hyperref[sol:g08]{G8}；附录\ref{app:stats} & 开放分析：扩题、重复与分层针对不同变异；重复与估计方法不增加独立题数。 \\
\hyperref[ex:ch04-5]{4-5} & 第\ref{sec:ch04-framing}节；\hyperref[sol:g15]{G15} & 开放设计：给旧目标之外的可观察结果，如后续重开或最终状态；说明反馈如何进入决策。 \\
\hyperref[ex:ch05-1]{5-1} & 表\ref{tab:ch05-forms}；\hyperref[sol:g18]{G18} & 开放观察：给实际条目、分类依据和可见失败，不以占比推断因果作用。 \\
\hyperref[ex:ch05-2]{5-2} & 式\eqref{eq:ch05-link}；\hyperref[sol:g09]{G9} & 开放反例：共享词义可以掩盖不同适用条件；列出排名及错误采用链。 \\
\hyperref[ex:ch05-3]{5-3} & 第\ref{ch05}章；\hyperref[sol:g16]{G16} & 开放分析：条件文本不等于物理删除、停止读取或撤销已传播内容；可给其他具体反例。 \\
\hyperref[ex:ch05-4]{5-4} & \hyperref[sol:g09]{G9}/\hyperref[sol:g10]{G10}；表\ref{tab:ch15-reader} & 开放分析：内容、来源质量和接触任务范围均可改变；先固定同源材料再解释写法差。 \\
\hyperref[ex:ch06-1]{6-1} & \hyperref[sol:g11]{G11}；表\ref{tab:ch06-who} & 开放判定：分别登记写、留、取、删；新增策略训练、反馈、运行与维护成本。 \\
\hyperref[ex:ch06-2]{6-2} & 式\eqref{eq:ch06-amac}；\hyperref[sol:g11]{G11} & 开放权衡：比较规则可复核性、模型语义判断与费用；不得预定性能必然上升。 \\
\hyperref[ex:ch06-3]{6-3} & 第\ref{ch06}章Mem-$\alpha$；\hyperref[sol:g15]{G15} & 开放设计：评价冲突发现与正确处置，并构造为得分而删除有效经验等投机路径。 \\
\hyperref[ex:ch06-4]{6-4} & 第\ref{ch06}章D-MEM；第\ref{ch13}章 & 开放设计：至少报告触发频率、结果质量与完整资源，并定义请求、尝试、发布的分母。 \\
\hyperref[ex:ch07-1]{7-1} & 第\ref{ch07}章更新流程；\hyperref[sol:g02]{G2} & 开放观察：比较版本差，区分删去、改写与被摘要吸收的旧内容。 \\
\hyperref[ex:ch07-2]{7-2} & 表\ref{tab:ch07-cost}；\hyperref[sol:g14]{G14} & 开放设计：交叉改变重写方式与候选评估调用量，统一任务、停止规则和计时范围。若存在交互，改变一个因素的节省量取决于另一个因素的取值。 \\
\hyperref[ex:ch07-3]{7-3} & 第\ref{ch07}章；\hyperref[sol:g07]{G7} & 开放设计：固定条件下干预条目是否可读，记录实际检索、输入及最终判定。 \\
\hyperref[ex:ch07-4]{7-4} & 表\ref{tab:ch15-reader}；\hyperref[sol:g10]{G10} & 开放设计：冻结同一上下文资产，以新读取者自身静态基线作对应比较。 \\
\hyperref[ex:ch08-1]{8-1} & 第\ref{ch08}章技能元数据；\hyperref[sol:g18]{G18} & 开放观察：先定义误用和暴露次数；库大小与任务难度同时变化时不能直接归因。 \\
\hyperref[ex:ch08-2]{8-2} & 表\ref{tab:ch08-voyager}；\hyperref[sol:g06]{G6} & 已有相同迭代上限下的有无对照；总资源比较还需来源构建、实际尝试及各阶段费用，固定可见反馈与选择机会。 \\
\hyperref[ex:ch08-3]{8-3} & 图\ref{fig:ch08-kstar}；\hyperref[sol:g12]{G12}/\hyperref[sol:g14]{G14} & 按图中假设计算：第4次平价，第5次严格节省。右图第6次加收维护费后技能较贵，第8次再次平价。若从第$m$次加收$M$，自该次起改比较$12+2k+M$与$5k$。 \\
\hyperref[ex:ch08-4]{8-4} & 第\ref{ch08}章；\hyperref[sol:g16]{G16} & 开放设计：预定技能集合、接口前后条件和静默失效标准，保留全部技能分母。 \\
\hyperref[ex:ch09-1]{9-1} & 图\ref{fig:ch09-workflow}；\hyperref[sol:g13]{G13} & 开放计算：直出时计生成求解与评审，转入修正时再加修正调用及等待。误放可能交付错解，误拒可能改坏正确解；净收益还需初答正确率、各分支修正效果和费用。 \\
\hyperref[ex:ch09-2]{9-2} & 图\ref{fig:ch09-workflow}；式\eqref{eq:ch09-latency} & 条件分析：总请求可能增加；时延由调度与队列决定，相关性由生成机制决定，无通用不变量。 \\
\hyperref[ex:ch09-3]{9-3} & 第\ref{ch09}章archive；\hyperref[sol:g18]{G18} & 开放设计：来源版本、评价范围与选择规则是最小起点，另列恢复依赖及跨搜索采用条件。 \\
\hyperref[ex:ch09-4]{9-4} & 第\ref{ch09}章预算；\hyperref[sol:g06]{G6} & 开放设计：固定共同反馈和总资源，预定扁平选择器；预算单位应匹配实际部署约束。 \\
\hyperref[ex:ch10-1]{10-1} & 表\ref{tab:ch10-seal}；\hyperref[sol:g06]{G6} & 开放分析：100是筛选范围内Oracle值；分别固定编辑内容与适配过程才能区分瓶颈。 \\
\hyperref[ex:ch10-2]{10-2} & 表\ref{tab:ch10-seal}、\ref{tab:ch10-cpt} & 单段使用 LoRA，多段使用全参数训练，较大语料档的评价范围也改变。停止规则须预定终评、训练量、真实数据比例和多次查看的处理。 \\
\hyperref[ex:ch10-3]{10-3} & 第\ref{sec:ch10-gate}节；\hyperref[sol:g16]{G16} & 开放判定：记忆条目与适配器分别绑定基座、读取配置及恢复快照；文件恢复不等于效果撤销。 \\
\hyperref[ex:ch10-4]{10-4} & 第\ref{sec:ch10-gradient}节；式\eqref{eq:ch10-normalization} & 开放反例：先说明等价指更新方向还是参数步，再看奖励、步数或筛选条件的改变是否保留原目标。 \\
\hyperref[ex:ch10-5]{10-5} & 第\ref{ch10}章资源；\hyperref[sol:g14]{G14} & 计算：串行且无额外开销时3000--4500秒，即50--75分钟；不能据此推算外部经验产量。 \\
\hyperref[ex:ch11-1]{11-1} & 第\ref{ch11}章章首；表\ref{tab:ch11-filter} & 冻结合成器可能使题目难度与求解器错配，去掉惩罚可能增加重复题。分别跟踪题目难度与改写重复率，保留共同的终评。 \\
\hyperref[ex:ch11-2]{11-2} & 第\ref{sec:ch11-filter}节；\hyperref[sol:g06]{G6} & 开放设计：交叉操纵合成器起点与可更新性，保持求解器、任务范围和资源一致。 \\
\hyperref[ex:ch11-3]{11-3} & 第\ref{sec:ch11-consensus}节 & 计算与设计。二元独立设定下，多数一致率期望为\calcnum{0.623046875}，正确率期望为\calcnum{0.5}。可行性、答案校验与难度分别评价，共同错误需独立依据。 \\
\hyperref[ex:ch11-4]{11-4} & 第\ref{sec:ch11-reward}节；\hyperref[sol:g15]{G15} & 开放定义。预定外部任务目标与可接受差异，再检查行为是否利用评价缺口抬高代理奖励。格式改善本身有价值时，需按该目标判断。 \\
\hyperref[ex:ch11-5]{11-5} & 第\ref{ch11}章；\hyperref[sol:g06]{G6}/\hyperref[sol:g15]{G15} & 开放诊断：选择偏差、训练重复、分布变易都可抬分，分别用封存题、去重和分布记录检查。 \\
\hyperref[ex:ch12-1]{12-1} & 表\ref{tab:ch12-sica}；\hyperref[sol:g04]{G4}/\hyperref[sol:g15]{G15} & 开放审计：代码共处只表明权限需查；实际改评测与反复选候选分别影响测量和独立性。 \\
\hyperref[ex:ch12-2]{12-2} & 第\ref{sec:ch12-archive}节；\hyperref[sol:g13]{G13} & 开放设计：缺编辑能力的高分终端与无效分支都可能入档，按实际选择与功能后代记录区分。 \\
\hyperref[ex:ch12-3]{12-3} & 表\ref{tab:ch12-imp}、\ref{tab:ch12-components} & 开放构造：同时列出起点、选定后代与差值；二者同幅提高时增量不变。修改器贡献还需在匹配起点、预算和组件安排后比较。 \\
\hyperref[ex:ch12-4]{12-4} & 第\ref{sec:ch12-vex}节；\hyperref[sol:g04]{G4} & 开放设计：新旧规则交叉复评或独立外部目标，保留规则版本与桥接费用。 \\
\hyperref[ex:ch12-5]{12-5} & 第\ref{sec:ch12-transfer}节；\hyperref[sol:g17]{G17} & 开放分析：大于零比较方法与自身起点；更可迁移比较方法间相同目标范围下的迁移增量。 \\
\hyperref[ex:ch13-1]{13-1} & 表\ref{tab:ch13-timing}；\hyperref[sol:g14]{G14} & 可把请求数除以任务数、尝试数除以请求数、发布数除以尝试数。在每次请求至多尝试一次、每次尝试至多发布一次时，三项分母均非零，其乘积就是每题发布数。若无请求，仍可报告每题发布数为零，后两项无分母。 \\
\hyperref[ex:ch13-2]{13-2} & \S\ref{sec:ch13-signals}；\hyperref[sol:g11]{G11} & 开放设计：观察信号分布、重复材料与真正新材料的区分，再验证其对目标的预测作用。 \\
\hyperref[ex:ch13-3]{13-3} & 第\ref{sec:ch13-costs}节；\hyperref[sol:g14]{G14} & 开放计算：按更新、验证、存储与检索、推理分别记每轮新增费用。若每轮扫描全部旧条目且各条处理费相同，前$n$轮的扫描总量随$n^2$增长，单轮随$n$增长。 \\
\hyperref[ex:ch13-4]{13-4} & \S\ref{sec:ch13-trace-example}；\hyperref[sol:g15]{G15} & 开放诊断：稳定、阈值过高与信号失灵均相容；补记原始信号、候选和监测运行状态。 \\
\hyperref[ex:ch14-1]{14-1} & \S\ref{sec:ch14-provenance-chain}；\hyperref[sol:g15]{G15} & 开放审计：依次查原标签、察觉事件、写入与读取轨迹；以原执行输出检验不一致。 \\
\hyperref[ex:ch14-2]{14-2} & \S\ref{sec:ch14-rail}；\hyperref[sol:g16]{G16} & 开放分析：哈希检验未改写，不能独自证明事实真实；摘要重算另需来源与转换版本。 \\
\hyperref[ex:ch14-3]{14-3} & 表\ref{tab:ch14-collapse}；\hyperref[sol:g15]{G15} & 开放设计：在固定任务上评价表现，独立审查接纳分布；类别熵可能掩盖类内同质化。 \\
\hyperref[ex:ch14-4]{14-4} & 第\ref{ch14}章框架锁定；\hyperref[sol:g15]{G15} & 开放设计：加入与实际目标相连的独立观察，说明其可见性和修正规则的权限。 \\
\hyperref[ex:ch14-5]{14-5} & 附录\ref{app:stats}；式\ref{eq:ch14-consistency}；\hyperref[sol:g05]{G5} & 开放设计：限定污染、任务和效应方向；精度需半宽，功效还需真值、显著性与错误控制；多次全成功与至少一次成功须分别计算。 \\
\hyperref[ex:ch15-1]{15-1} & 表\ref{tab:ch15-reader}；式\eqref{eq:ch15-reader-interaction} & 计算：两差分别为 $-0.3250$、$+0.1000$，后者是目标读取者内的增量。另在每个读取者内取领域拟合减求解器拟合，再用换读差减同读差，得交互点值 $0.0625$。 \\
\hyperref[ex:ch15-2]{15-2} & 表\ref{tab:ch05-e03a}；\hyperref[sol:g09]{G9} & 审计：拆分A缺少能解释两次判定差异的记录。暂停该静态基线、三种记忆臂相对它的差值与区间，以及依赖它的成本效率解释。三种记忆臂之间仍可直接比较。 \\
\hyperref[ex:ch15-3]{15-3} & 表\ref{tab:ch15-validity}；\hyperref[sol:g10]{G10} & 开放设计：预定共同基线、任务与场景、过程合规检查和成本。样本量须对应主指标及场景内相关性；若主张迁移效果不变，应事先规定等价范围。 \\
\hyperref[ex:ch15-4]{15-4} & 表\ref{tab:appF-cross} & 固定总预算时，少做题会缩小任务覆盖，减少重复会降低逐题估计的精度，缩短输出则会改变允许求解器使用的资源。 \\
\hyperref[ex:ch15-5]{15-5} & 第\ref{ch05}章双字段；\hyperref[sol:g09]{G9} & 开放分析：题级收益点估计反号不证明真实效应反号；双字段的唯一必要性还需替代设计对照。 \\
\hyperref[ex:ch16-1]{16-1} & \hyperref[sol:g05]{G5}；表\ref{tab:ch16-power} & 统计解释：半宽式回答精度问题；功效设计另需零假设、备择效应、显著性水平和目标功效，达到精度题数仍不能修复设计、单位与覆盖。 \\
\hyperref[ex:ch16-2]{16-2} & \hyperref[sol:g17]{G17}；表\ref{tab:ch16-verdicts} & 开放诊断：P1问预算分配，P3问形式与容量交互，P5问时机；分别需要匹配重试、交叉对照和任务流比较。 \\
\hyperref[ex:ch16-3]{16-3} & \hyperref[sol:g09]{G9}/\hyperref[sol:g10]{G10}；第\ref{ch05}章 & 开放设计：形式乘容量，固定来源与检索预算，预定独立族与比较范围。 \\
\hyperref[ex:ch16-4]{16-4} & 第\ref{sec:ch14-claims}节 & $D$ 为旧分减新分。区间含零时方向未定；若范围内某次更新的有效配对比较支持 $D<0$，该例就与普遍退化的主张冲突。 \\
\hyperref[ex:ch16-5]{16-5} & \hyperref[sol:g17]{G17}；表\ref{tab:ch15-reader} & 审计：补齐的是一臂记录；跨日期服务身份、少数读取者与未换harness的范围仍需限定。 \\
\end{longtable}
\end{small}
